%!TEX program = pdflatex
\documentclass[11pt,a4paper]{article}

% ===== PACKAGES =====
\usepackage[T1]{fontenc}
\usepackage[utf8]{inputenc}
\usepackage{lmodern}
\usepackage[margin=1in]{geometry}
\usepackage{amsmath,amssymb,amsthm}
\usepackage{mathtools}
\usepackage{physics}
\usepackage{tikz}
\usepackage{pgfplots}
\pgfplotsset{compat=1.18}
\usetikzlibrary{arrows.meta,positioning,decorations.markings,decorations.pathmorphing,calc,patterns,shapes.geometric}

% Define photon line style for Feynman diagrams
\tikzset{
    photon/.style={decorate, decoration={snake, amplitude=1.5pt, segment length=5pt}}
}
\usepackage{booktabs}
\usepackage{array}
\usepackage{xcolor}
\usepackage{float}
\usepackage[hidelinks]{hyperref}
\usepackage{microtype}
\usepackage{enumitem}
\usepackage{pifont}
\usepackage{siunitx}
\usepackage{tcolorbox}
\tcbuselibrary{theorems,skins,breakable}

% ===== COLORS =====
\definecolor{axiomcolor}{RGB}{0,80,160}
\definecolor{theoremcolor}{RGB}{0,120,60}
\definecolor{corollarycolor}{RGB}{120,60,0}
\definecolor{definitioncolor}{RGB}{100,0,100}
\definecolor{topologyblue}{RGB}{50,100,180}
\definecolor{geometrygreen}{RGB}{50,150,80}
\definecolor{quantumorange}{RGB}{220,120,40}
\definecolor{observablepurple}{RGB}{140,60,160}

% ===== THEOREM ENVIRONMENTS =====
\newtcbtheorem[number within=section]{axiom}{Axiom}{
    colback=axiomcolor!5,
    colframe=axiomcolor,
    fonttitle=\bfseries,
    separator sign={.},
    description delimiters parenthesis
}{ax}

\newtcbtheorem[number within=section]{theorem}{Theorem}{
    colback=theoremcolor!5,
    colframe=theoremcolor,
    fonttitle=\bfseries,
    separator sign={.},
    description delimiters parenthesis
}{thm}

\newtcbtheorem[number within=section]{corollary}{Corollary}{
    colback=corollarycolor!5,
    colframe=corollarycolor,
    fonttitle=\bfseries,
    separator sign={.},
    description delimiters parenthesis
}{cor}

\newtcbtheorem[number within=section]{lemma}{Lemma}{
    colback=theoremcolor!5,
    colframe=theoremcolor,
    fonttitle=\bfseries,
    separator sign={.},
    description delimiters parenthesis
}{lem}

\newtcbtheorem[number within=section]{definition}{Definition}{
    colback=definitioncolor!5,
    colframe=definitioncolor,
    fonttitle=\bfseries,
    separator sign={.},
    description delimiters parenthesis
}{def}

\newenvironment{remark}{\par\smallskip\noindent\textbf{Remark.}}{\par\smallskip}

% ===== COMMANDS =====
\newcommand{\Mpl}{\bar{M}_{\mathrm{Pl}}}
\newcommand{\cthree}{c_3}
\newcommand{\phiz}{\varphi_0}
\newcommand{\phitree}{\varphi_{\mathrm{tree}}}
\newcommand{\deltatop}{\delta_{\mathrm{top}}}
\newcommand{\bone}{b_1}
\newcommand{\gagg}{g_{a\gamma\gamma}}

% ===== TITLE =====
\title{%
    \textbf{Topological Fixed Point Theory (TFPT)}\\[0.5em]
    \Large A Parameter-Free Derivation of the Fine-Structure Constant\\
    from Topology, Geometry, and Quantum Consistency
}
\author{%
    Stefan Hamann \and Alessandro Rizzo
}
\date{Version 2.2 --- 21 January 2026}

